\documentclass[11pt,a4paper]{article}
\usepackage{fontspec}
\usepackage{geometry}
\usepackage{amsmath,amssymb,mathtools}
\usepackage{booktabs}
\usepackage{enumitem}
\usepackage{xcolor}
\usepackage{titlesec}
\usepackage{fancyhdr}
\usepackage[unicode=true]{hyperref}
\usepackage[most]{tcolorbox}
\usepackage{lastpage}
\usepackage{microtype}
\geometry{top=22mm,bottom=24mm,left=24mm,right=24mm,headheight=14pt}
\definecolor{ink}{HTML}{172033}\definecolor{accent}{HTML}{315A8A}\definecolor{accentlight}{HTML}{EAF1F8}\definecolor{rulegray}{HTML}{D7DEE8}\definecolor{muted}{HTML}{5C6678}
\setmainfont{DejaVu Serif}\setsansfont{DejaVu Sans}\setmonofont{DejaVu Sans Mono}
\hypersetup{colorlinks=true,linkcolor=accent,urlcolor=accent,citecolor=accent,pdftitle={Simetrías de semiperíodo para sucesiones módulo 3},pdfsubject={Formalización matemática y auditoría algorítmica},pdfkeywords={recurrencia, módulo 3, antiperiodicidad, Fourier, matriz compañera}}
\pagestyle{fancy}\fancyhf{}\fancyhead[L]{\small\color{muted}Simetrías de semiperíodo módulo 3}\fancyhead[R]{\small\color{muted}Formalización y auditoría}\fancyfoot[C]{\small\color{muted}\thepage\,/\,\pageref*{LastPage}}\renewcommand{\headrulewidth}{0.3pt}\renewcommand{\footrulewidth}{0pt}
\titleformat{\section}{\Large\bfseries\color{ink}}{\thesection.}{0.6em}{}[\vspace{0.2em}\color{rulegray}\titlerule]\titleformat{\subsection}{\large\bfseries\color{accent}}{\thesubsection}{0.6em}{}\titlespacing*{\section}{0pt}{2.2ex plus .5ex minus .2ex}{1.3ex}\titlespacing*{\subsection}{0pt}{1.7ex plus .4ex minus .2ex}{0.8ex}
\setlist[itemize]{leftmargin=1.6em,itemsep=0.25em,topsep=0.4em}\setlist[enumerate]{leftmargin=1.8em,itemsep=0.25em,topsep=0.4em}\setlength{\parindent}{0pt}\setlength{\parskip}{0.55em}\allowdisplaybreaks
\newtcolorbox{summarybox}{colback=accentlight,colframe=accent,boxrule=0.6pt,arc=2mm,left=4mm,right=4mm,top=3mm,bottom=3mm,before skip=1em,after skip=1em}\newcommand{\F}{\mathbb F}\newcommand{\Z}{\mathbb Z}\newcommand{\concat}{\mathbin{\Vert}}
\begin{document}
\begin{titlepage}\thispagestyle{empty}\vspace*{22mm}{\color{accent}\rule{\textwidth}{1.2pt}}\\[15mm]{\fontsize{27}{33}\selectfont\bfseries\color{ink}Simetrías de semiperíodo para sucesiones módulo 3\par}\vspace{9mm}{\Large\color{accent}Formalización matemática y auditoría algorítmica\par}\vspace{13mm}\begin{summarybox}Esta nota distingue rigurosamente tres fenómenos que suelen confundirse: la negación de ciclos, la antiperiodicidad de semiperíodo y la inversión. Demuestra el único desplazamiento opuesto de una palabra primitiva no nula, ofrece caracterizaciones combinatorias, espectrales y matriciales y relaciona estos resultados con un analizador exhaustivo de recurrencias modulares.\end{summarybox}\vfill{\large\color{muted}Repositorio: \texttt{Modular-recurrence-symetry-analyzer-}\par}{\large\color{muted}Edición PDF consolidada - 5 de agosto de 2026\par}\vspace{12mm}{\color{accent}\rule{\textwidth}{1.2pt}}\end{titlepage}
\renewcommand{\contentsname}{Índice}\tableofcontents\clearpage
\section{Objeto}
Esta nota formaliza los fenómenos descritos en el repositorio \href{https://github.com/59200/k-bonacci-pisano-finder}{\texttt{59200/k-bonacci-pisano-finder}} mediante las expresiones informales «períodos complementarios a tres» y «autocomplementarios cuando se cortan en dos».

El marco invariante es el de las palabras periódicas sobre el cuerpo \(\F_3=\{0,1,2\}\), con involución aditiva \(\nu(x)=-x\pmod3\), es decir, \(\nu(0)=0\), \(\nu(1)=2\), \(\nu(2)=1\). Por tanto, «complemento a 3» debe sustituirse por \textbf{opuesto aditivo módulo 3} o \textbf{complemento por negación módulo 3}. La suma decimal de los representantes no siempre vale 3, pues 0 se transforma en 0.
\section{Tres propiedades distintas}
Sea \(w=(w_0,\ldots,w_{T-1})\in\F_3^T\) una palabra de período primitivo \(T\).
\subsection{Par de períodos opuestos}
Dos períodos \(w\) y \(v\) forman un par opuesto si \(v_n=-w_n\pmod3\) para todo \(n\). Pueden pertenecer a dos ciclos distintos.
\subsection{Antiperiodicidad de semiperíodo}
La palabra \(w\) es \textbf{antiperiódica de semiperíodo} si \(T=2h\) y \(w_{n+h}=-w_n\pmod3\) para todo \(n\). Entonces \(w=H\concat(-H)\), con \(H=(w_0,\ldots,w_{h-1})\). Ejemplos: \(01120221=0112\concat0221\), \(011022=011\concat022\), \(01220211=0122\concat0211\).
\subsection{Inversión de la segunda mitad}
La operación \(\mathcal R(a_0,\ldots,a_{h-1})=(a_{h-1},\ldots,a_0)\) es independiente de la negación. En el ejemplo de Fibonacci, \(\mathcal R(0221)=1220\). Así, \(1220\) es la \textbf{palabra antipodal invertida} de \(0112\), y no la segunda mitad en sí.
\section{Teorema del único desplazamiento opuesto}
\begin{summarybox}\textbf{Teorema.} Sea \(w\) una palabra primitiva no nula de período \(T\). Si existe un desplazamiento \(r\) tal que \(w_{n+r}=-w_n\) para todo \(n\), entonces \(T\) es par y \(r\equiv T/2\pmod T\).\end{summarybox}
\subsection*{Demostración}
Aplicando dos veces el desplazamiento se obtiene \(w_{n+2r}=w_n\). La primitividad implica \(T\mid2r\). El desplazamiento \(r\) no es nulo módulo \(T\), porque una palabra no nula sobre \(\F_3\) no puede satisfacer \(w=-w\). El único elemento no trivial de orden 2 del grupo cíclico \(\Z/T\Z\) existe solo cuando \(T\) es par y es \(T/2\). \hfill\(\square\)
\subsection*{Consecuencia}
Un período primitivo no nulo es, por tanto: \begin{enumerate}\item enviado por la negación a un ciclo distinto;\item o invariante por negación, caso en el que es necesariamente antiperiódico de semiperíodo.\end{enumerate} Esta dicotomía elimina la confusión entre el par \(00111201,00222102\) y períodos intrínsecamente antiperiódicos como \(011022\).
\section{Consecuencias combinatorias}
Si \(w=H\concat(-H)\), el número de 1 coincide con el número de 2, el número de 0 es par y \(\sum_{n=0}^{T-1}w_n=0\) en \(\F_3\). Estas condiciones son necesarias, pero no suficientes. Con el levantamiento centrado \(\widetilde0=0\), \(\widetilde1=1\), \(\widetilde2=-1\), el polinomio generador factoriza como \[W(X)=\sum_{n=0}^{2h-1}\widetilde w_nX^n=(1-X^h)\sum_{n=0}^{h-1}\widetilde w_nX^n.\]
\section{Caracterización de Fourier}
Sea \(\widehat w(k)=\sum_{n=0}^{2h-1}\widetilde w_ne^{-2\pi i kn/(2h)}\). Entonces \(w_{n+h}=-w_n\) si y solo si \(\widehat w(k)=0\) para todo índice par \(k\). En el sentido directo aparece el factor \(1-(-1)^k\), que se anula para \(k\) par. Recíprocamente, si todos los modos pares son nulos, la palabra pertenece al subespacio generado por los modos impares y el desplazamiento de \(h\) actúa como \(-I\). Es un diagnóstico espectral exacto, no una simple comparación visual.
\section{Sucesiones lineales módulo 3}
Consideremos la recurrencia de orden \(k\), \[u_{n+k}=c_0u_n+\cdots+c_{k-1}u_{n+k-1}\pmod3.\] Su estado es \(s_n=(u_n,\ldots,u_{n+k-1})^{\mathsf T}\) y evoluciona según \(s_{n+1}=Ms_n\), donde \(M\) es la matriz compañera.
\subsection{Periodicidad pura}
El determinante de \(M\) es, salvo signo, \(c_0\). Sobre \(\F_3\), \(M\) es invertible si y solo si \(c_0\neq0\). En ese caso toda órbita de estados es puramente periódica. Si \(c_0=0\), puede aparecer un preperíodo que debe distinguirse del ciclo.
\subsection{Negación de los ciclos}
La linealidad da \(M(-s)=-M(s)\). La negación envía cada ciclo a otro de la misma longitud. Para un ciclo no nulo solo hay dos casos: dos ciclos distintos intercambiados por negación, o un único ciclo invariante, necesariamente antiperiódico de semiperíodo. El ciclo nulo es la única excepción degenerada: queda fijado punto por punto y no define un antiperíodo primitivo no trivial.
\subsection{Criterio para una semilla}
Para una semilla \(s_0\), la relación \(M^hs_0=-s_0\) implica \(u_{n+h}=-u_n\) para toda salida lineal del estado. Cuando \(s_0\neq0\) y el ciclo es primitivo, proporciona el antiperíodo de semiperíodo descrito antes.
\subsection{Criterio uniforme}
Todas las semillas satisfacen la misma relación de antiperíodo \(h\) si y solo si \(M^h=-I\). Equivalentemente, \(\mu_M(X)\mid X^h+1\) en \(\F_3[X]\); de ahí se sigue \(M^{2h}=I\).
\section{Ejemplo de Fibonacci}
Para \(M=\begin{pmatrix}0&1\\1&1\end{pmatrix}\) sobre \(\F_3\), se obtiene \(M^4=-I\) y \(M^8=I\). La semilla \((0,1)\) genera \(01120221=0112\concat(-0112)\). La segunda mitad invertida es \(\mathcal R(0221)=1220\). La coincidencia decimal \(1220=L_{14}+F_{14}=2F_{15}\) es una identidad adicional de codificación, no una consecuencia general de la antiperiodicidad.
\section{Generalización a sucesiones arbitrarias módulo 3}
No se necesita ninguna hipótesis de recurrencia para analizar una palabra periódica dada. Para cada período primitivo se pueden calcular su ciclo de rotaciones, su opuesto aditivo, sus simetrías afines \(w_{n+r}=aw_n+b\), sus simetrías de inversión \(w_{r-n}=aw_n+b\), su posible antiperiodicidad y sus modos de Fourier pares e impares. La teoría matricial solo interviene cuando se dispone de una recurrencia lineal declarada.
\section{Extensión a un módulo general}
Sobre \(\Z/m\Z\), la involución natural sigue siendo \(x\mapsto-x\), y \(w_{n+h}=-w_n\pmod m\) conserva su sentido. Para \(m=2\), la negación es la identidad y la antiperiodicidad coincide con la periodicidad ordinaria; para \(m>2\), la distinción es no trivial. «Complemento a \(m\)» no es una definición algebraica invariante; el término correcto es «opuesto aditivo módulo \(m\)».
\section{Auditoría algorítmica}
El período de una recurrencia de orden \(k\) módulo \(m\) debe detectarse en el espacio finito de estados \((\Z/m\Z)^k\), no buscando un bloque repetido en un prefijo arbitrariamente largo. El script usa codificación compacta en base \(m\), el algoritmo de Brent para una semilla en tiempo \(O(\mu+\lambda)\) y memoria \(O(1)\), descomposición exacta del grafo funcional en tiempo \(O(m^k)\), paralelización entre familias, clasificación separada de ciclos opuestos y antiperiódicos y una comprobación matricial de \(M^h=-I\). No requiere Mathematica ni envoltorios externos.
\section{Resultados reproducibles del barrido}
El barrido de las quince familias nombradas recupera, para la familia tritetranacci, pares opuestos distintos de longitudes 24 y 8 y los períodos antiperiódicos \(011022\), \(01220211\) y \(12\). Un segundo barrido considera los \(119\) vectores binarios no nulos de órdenes 2 a 6: \(\sum_{k=2}^{6}(2^k-1)=119\). Se obtienen 13 familias con criterio global \(M^h=-I\), 88 con al menos un ciclo antiperiódico no trivial y 96 con al menos un par distinto de ciclos opuestos. Los resultados son reproducibles a partir de los archivos JSON y CSV del repositorio.
\section*{Conclusión}\addcontentsline{toc}{section}{Conclusión}
La negación módulo 3 actúa como una involución estructural sobre las palabras periódicas y los ciclos de recurrencias lineales. Para una palabra primitiva no nula, la invariancia solo puede producirse mediante el único desplazamiento de semiperíodo. La misma propiedad se reconoce por la factorización del polinomio generador, la anulación de los modos de Fourier pares o, en el marco lineal, la relación matricial \(M^h=-I\). El análisis algorítmico separa así las simetrías intrínsecas de las simples asociaciones entre ciclos opuestos.
\end{document}